% Appendix.
%
% The evidence a forgetting benchmark owes but that does not belong in the
% results: that the task being forgotten was learned, that the budget sits
% somewhere defensible on the convergence curve, and that the numbers can be
% reproduced. Tables under tables/ are generated.

\appendix

\section{Was there anything to forget?}
\label{sec:appendix:quality}

\begin{table}[h]
  \centering
  % Generated by experiments/export_tables.py -- do not edit by hand.
% Regenerate after any change to results/ or to the aggregation.

\begin{tabular}{llrrr}
\toprule
Family & Level & Held-out recon. & NLL at init & NLL after A \\
\midrule
MiniGrid & min & 0.48 & 37.19 & 14.89 \\
 & med & 0.84 & 38.65 & 13.78 \\
 & max & 42.58 & 37.45 & 15.65 \\
\addlinespace
Gymnasium & min & 26.49 & 41.30 & 12.42 \\
 & med & 26.15 & 41.30 & 12.42 \\
 & max & 26.82 & 41.50 & 12.59 \\
\addlinespace
DMControl & min & 14.74 & 29.61 & 10.68 \\
 & med & 14.54 & 29.63 & 11.73 \\
 & max & 14.54 & 29.63 & 11.73 \\
\bottomrule
\end{tabular}

  \caption{Task-A quality per cell, measured before any task switch and pooled
  over methods --- at that point in the protocol every method has trained on
  task~A and nothing else. \emph{Held-out recon.} is squared error summed over
  the $3 \times 64 \times 64$ frame on episodes never trained on. The two NLL
  columns are the transition component's one-step negative log-likelihood on
  $D_A$ at initialisation and after training, so their difference is what
  training bought.}
  \label{tab:quality}
\end{table}

A forgetting benchmark that has not shown its models learned the first task is
reporting the distance between two bad models. Table~\ref{tab:quality} is that
check, for every cell. Training on task~A moves the one-step NLL on $D_A$ from
the high thirties or low forties to between $10.68$ and $15.65$ in every cell,
and the held-out reconstruction reaches values that differ by two orders of
magnitude between the visually simplest cells and the hardest --- which is the
task-difficulty signal Section~\ref{sec:results:predictors} uses, seen at the
first task instead of the second.

The one cell that deserves comment is MiniGrid at maximum distance, whose
held-out reconstruction of task~A is $42.58$ where the other two MiniGrid cells
are below $1$. Its task~A is FourRooms, which is the hardest environment in that
family to fit; the model still improves its NLL on $D_A$ by more than a factor
of two, so there is something learned there to lose, but the pixel-scale
comparisons in that row are between larger numbers than elsewhere.

\section{Where the budget sits on the convergence curve}
\label{sec:appendix:budget}

The protocol trains 5000 gradient updates per task. That number was chosen from
a scaling probe on MiniGrid rather than by default: a single run evaluated at
1000, 2000, 5000 and 10000 updates takes held-out task-A reconstruction from
$6.49$ to $3.01$, $1.61$ and $1.37$. The step from 1000 to 5000 removes three
quarters of the error; the step from 5000 to 10000 removes a further 15\%.

We report that as locating the budget past the steep part of the curve and
short of its floor, which is what a forgetting benchmark needs --- a model that
has not converged has less to lose, and one that has fully converged is not the
regime anybody trains in either. It does not license reading the absolute
magnitudes as budget-independent, and Section~\ref{sec:results:axis} is the one
place where we tested the direction of a result against a change of budget
rather than assuming it.

These four numbers come from a separate probe rather than from the grid, which
is why they do not appear in a generated table.

\section{Reproducibility}
\label{sec:appendix:reproducibility}

Every result in this paper is produced by two scripts that train nothing: one
prints the console summary and one emits the paper's tables in \LaTeX{} by
importing the first. Every table in the paper carries a header saying it was
generated. This is not tidiness --- the previous version of this work
transcribed its tables by hand and they drifted from the runs they claimed to
describe, and three different sets of protocol values circulated for the same
experiment.

Runs are reproducible from their seed, and the part that took longest to get
right was not the framework's global RNG but each environment's own: Gymnasium's
action space and dm\_control's \texttt{RandomState} are not reached by seeding
NumPy, and until they were seeded the five seeds of a cell trained on five
different datasets and were not replicates. All instrumentation is wrapped so
that measuring cannot perturb a run, which is load-bearing because UG-MTM's gate
keeps MC-dropout active at evaluation time by design and so consumes the random
stream. We reproduced a cell bit for bit eight days after it first ran,
in a different process, after the runner had been rewritten in between.

Each result file records the protocol block it was produced under and the task
pair it was run on. The runner refuses to reuse a cached result that disagrees
with either, the loader refuses to aggregate a directory whose runs disagree,
and the table exporter refuses to emit from one. Those three checks share one
definition of what makes two results comparable.

\section{Where the measured distance fails}
\label{sec:appendix:dtrans}

\begin{figure}[h]
  \centering
  \includegraphics[width=\textwidth]{figures/forgetting_vs_distance}
  \caption{The same runs as Figure~\ref{fig:axis}, plotted against the measured
  distance $d_{\mathrm{trans}}$ instead of the labelled level. Compare the
  x-axis ticks: within each family the three levels sit almost on top of one
  another, and in DMControl two of the three labels collide. This is
  Section~\ref{sec:results:predictors}'s point made visually --- the measured
  distance separates families cleanly and does not resolve the levels inside
  one.}
  \label{fig:dtrans}
\end{figure}

Figure~\ref{fig:dtrans} is the figure this paper would have led with under its
earlier framing, when $d_{\mathrm{trans}}$ was believed to be the better axis.
We include it because what it shows now is the opposite of what it was drawn
for, and because the failure is easier to see than to state: the cells do not
spread out along the measured axis, and re-sorting them by it is what makes the
peak of Figure~\ref{fig:axis} invisible.
